% LaTeX Package: runcode 2025/01/17 v2.4
% 
% Copyright (C) 2020-2025 by Haim Bar and HaiYing Wang
% 
% This file may be distributed and/or modified under the
% conditions of the LaTeX Project Public License, either
% version 1.3c of this license or (at your option) any later
% version.  The latest version of this license is in:
% 
% http://www.latex-project.org/lppl.txt
% 
% and version 1.3c or later is part of all distributions of
% LaTeX version 2005/12/01 or later.

\documentclass{ltxdoc}
\usepackage[utf8]{inputenc}
\usepackage[T1]{fontenc}
\usepackage{graphicx}
\usepackage{longtable}
\usepackage{wrapfig}
\usepackage{rotating}
\usepackage[normalem]{ulem}
\usepackage{amsmath}
\usepackage{amssymb}
\usepackage{capt-of}
\usepackage{hyperref}
\usepackage{minted}
\author{Haim Bar and HaiYing Wang \\haim.bar@uconn.edu, haiying.wang@uconn.edu}
\date{\today}
\title{The \textbf{runcode} package}
\begin{document}

\maketitle
\begin{abstract}

\texttt{runcode} is a \LaTeX{} package that executes programming source codes (including
all command line tools) from \LaTeX{}, and embeds the results in the resulting pdf
file. Many programming languages can be easily used and any command-line
executable can be invoked when preparing the pdf file from a tex file. \texttt{runcode}
is also available on \href{https://ctan.org/pkg/runcode}{CTAN}.

It is recommended to use this package in the server mode together with the
\href{https://www.python.org/}{Python} \href{https://pypi.org/project/talk2stat/}{talk2stat} package. Currently, the server mode supports \href{https://julialang.org/}{Julia}, \href{https://www.mathworks.com/products/matlab.html}{MatLab},
\href{https://www.python.org/}{Python}, and \href{https://www.r-project.org/}{R}. More languages will be added.

For more details and usage examples and troubleshooting, refer to the
package’s github repository, at \url{https://github.com/Ossifragus/runcode}.

\end{abstract}

\section{Installation}
\label{sec:org7b41ef7}
You can simply put the runcode.sty file in the \LaTeX{} project folder.

The server mode requires the \href{https://pypi.org/project/talk2stat/}{talk2stat} package, and some chunk command (e.g.,
\texttt{\textbackslash{}showChunk}) requires the \href{https://pypi.org/project/advance-touch/}{advance-touch} package. To install them from the
command line, use:

\begin{verbatim}
pip3 install talk2stat advance-touch
\end{verbatim}

\textbf{Note}: \texttt{runcode} requires to enable the \texttt{shell-escape} option when
compiling a \LaTeX{} document.


\section{Usage}
\label{sec:org5034faf}
\subsection{Load the package:}
\label{sec:orge3286de}
\begin{minted}[]{latex}
\usepackage[options]{runcode}
\end{minted}

Available options are:

\begin{itemize}
\item \texttt{cache}: use cached results.

\item \texttt{fvextra}: use the \href{https://ctan.org/pkg/fvextra}{fvextra} package to show code.

\item \texttt{julia}: start server for \href{https://julialang.org/}{Julia} (requires \href{https://pypi.org/project/talk2stat/}{talk2stat}).

\item \texttt{listings}: use the \href{https://ctan.org/pkg/listings?lang=en}{listings} package to show code.

\item \texttt{matlab}: start server for \href{https://www.mathworks.com/products/matlab.html}{MatLab} (requires \href{https://pypi.org/project/talk2stat/}{talk2stat}).

\item \texttt{minted}: use the \href{https://ctan.org/pkg/minted}{minted} package to show code (requires the \href{https://pygments.org/}{pygments} package).
This is the default option.

\item \texttt{nominted}: use the \href{https://ctan.org/pkg/fvextra}{fvextra} package
instead of the \href{https://ctan.org/pkg/minted}{minted} package to show
code (this does not require the \href{https://pygments.org/}{pygments}
package, but it does not provide syntax highlights).

\item \texttt{nohup}: use the \texttt{nohup} command when starting a server. When using
the server-mode, some editors terminate all child processes after
\LaTeX{} compiling such as Emacs with Auctex. This option set the
variable notnohup to be false, and the server will not be terminated
by the parent process. \textbf{This option has to be declared before
declaring any language}, e.g., \texttt{[nohup, R]} works but \texttt{[R, nohup]}
does not work.

\item \texttt{python}: start server for \href{https://www.python.org/}{Python}
(requires \href{https://pypi.org/project/talk2stat/}{talk2stat}).

\item \texttt{run}: run source code.

\item \texttt{R}: start server for \href{https://www.r-project.org/}{R} (requires
\href{https://pypi.org/project/talk2stat/}{talk2stat}).

\item \texttt{stopserver}: stop the server(s) when the pdf compilation is done.
\end{itemize}

\textbf{Note}: If \href{https://ctan.org/pkg/minted}{minted} is used, the style of
the code block is controlled through the minted package,
\href{https://github.com/Ossifragus/runcode/blob/master/examples/MontyHall/MontyHall.tex\#L3-L4}{e.g.:}

\begin{minted}[]{latex}
\setminted[julia]{linenos, frame=single, bgcolor=bg, breaklines=true}
\setminted[R]{linenos, frame=single, bgcolor=lightgray, breaklines=true}
\end{minted}

Similarly, \href{https://ctan.org/pkg/fvextra}{fvextra} and \href{https://ctan.org/pkg/listings?lang=en}{listings} packages can be customized through the \texttt{\textbackslash{}fvset}
and \texttt{\textbackslash{}lstset} commands, respectively, e.g.: 

\begin{minted}[]{latex}
\fvset{fontsize=\small, linenos=true, frame=single}
\lstset{basicstyle=\large, frame=single}
\end{minted}

The outputs from executing codes are displayed in
\href{https://ctan.org/pkg/tcolorbox?lang=en}{tcolorbox}, so the style can
be customized with \texttt{\textbackslash{}tcbset},
\href{https://github.com/Ossifragus/runcode/blob/master/examples/MontyHall/MontyHall.tex\#L5}{e.g.:}

\begin{minted}[]{latex}
\tcbset{breakable,colback=red!5!white,colframe=red!75!black}
\end{minted}

\subsection{Basic commands:}
\label{sec:org496679c}
\begin{itemize}
\item \texttt{\textbackslash{}runExtCode\{Arg1\}\{Arg2\}\{Arg3\}[Arg4]} runs an external code.

\begin{itemize}
\item \texttt{Arg1} is the executable program.
\item \texttt{Arg2} is the source file name.
\item \texttt{Arg3} is the output file name (with an empty value, the counter \texttt{codeOutput} is
used).
\item \texttt{Arg4} controls whether to run the code. \texttt{Arg4} is optional with three possible
values.
\begin{itemize}
\item empty value or skipped: the value of the global Boolean variable \texttt{runcode}
is used;
\item \texttt{run}: the code will be executed;
\item \texttt{cache} (or anything else): use cached results; if the output file does not
exist, the cache option will be overridden and the code will run.
\end{itemize}
\end{itemize}

\item \texttt{\textbackslash{}showCode\{Arg1\}\{Arg2\}[Arg3][Arg4]} shows the source code, using
\href{https://ctan.org/pkg/minted}{minted} (requires \href{https://pygments.org/}{pygments}), \href{https://ctan.org/pkg/fvextra}{fvextra}, or \href{https://ctan.org/pkg/listings?lang=en}{listings}.

\begin{itemize}
\item \texttt{Arg1} is the programming language.
\item \texttt{Arg2} is the source file name.
\item \texttt{Arg3} is the first line to show (optional with a default value 1).
\item \texttt{Arg4} is the last line to show (optional with a default value of
the last line).
\end{itemize}

\item \texttt{\textbackslash{}includeOutput\{Arg1\}[Arg2]} is used to embed the output from executed
code.

\begin{itemize}
\item \texttt{Arg1} is the output file name, and it needs to have the same value
as that of \texttt{Arg3} in \texttt{\textbackslash{}runExtCode}. If an empty value is given to
\texttt{Arg1}, the counter \texttt{codeOutput} is used.
\item \texttt{Arg2} is optional and it controls the type of output with a default
value \texttt{vbox}.
\begin{itemize}
\item \texttt{vbox} (or skipped) = verbatim in a box;
\item \texttt{tex} = pure latex;
\item \texttt{inline} = embed result in text.
\end{itemize}
\end{itemize}

\item \texttt{\textbackslash{}inln\{Arg1\}\{Arg2\}[Arg3][Arg4]} is designed for simple calculations; it runs
one command (or a short batch) and displays the output within the
text.

\begin{itemize}
\item \texttt{Arg1} is the executable program or programming language.
\item \texttt{Arg2} is the source code.
\item \texttt{Arg3} is the output file name (optional);
\item \texttt{Arg4} is the output type and controls whether to run the code.
\begin{itemize}
\item \texttt{inline} (or skipped or with empty value): embed result in text;
\item \texttt{vbox}: verbatim in a box;
\item \texttt{cache} or appending \texttt{.cache} to the argument (e.g., \texttt{vbox.cache}) AND the
output file in \texttt{Arg3} exists: the cached result will be used and the code
will not run.
\end{itemize}
\end{itemize}

\item \texttt{\textbackslash{}showChunk\{Arg1\}\{Arg2\}[Arg3][Arg4]} prints a selected chunk from source
code. The chunk is identified in the source code by two strings that define
the beginning and end of the chunk. The default beginning is \texttt{label===<name>}
where \texttt{<name>} should be a unique user-defined chunk ID. The default end marker is
\texttt{===end}. In the code, these markers should appear after a comment character, so
that the code will run.

\begin{itemize}
\item \texttt{Arg1} is the programming language.
\item \texttt{Arg2} is the source file name.
\item \texttt{Arg3} is is the chunk identifier.
\item \texttt{Arg4} and \texttt{Arg5} are the beginning and end markings of a chunk (optional with
default \texttt{label==} and \texttt{===end}.
\end{itemize}
\end{itemize}



\subsection{Extended commands:}
\label{sec:orgf1638ce}
\begin{itemize}
\item \texttt{\textbackslash{}runCodeIncOut\{Arg1\}\{Arg2\}[Arg3][Arg4][Arg5]} runs an external code and
embeds the output. This is a combination of \texttt{\textbackslash{}runExtCode} and \texttt{\textbackslash{}includeOutput}.
\begin{itemize}
\item \texttt{Arg1} is the executable program.
\item \texttt{Arg2} is the source file name.
\item \texttt{Arg3} (optional) controls whether to run the code. Its functionality is the
same as that of \texttt{Arg4} of \texttt{\textbackslash{}runExtCode}.
\item \texttt{Arg4} (optional) is the output file name. Its functionality is the same as
that of \texttt{Arg3} of \texttt{\textbackslash{}runExtCode}.
\item \texttt{Arg5} (optional) controls the type of output with a default value \texttt{vbox}. Its
functionality is the same as that of \texttt{Arg3} of \texttt{\textbackslash{}includeOutput}.
\end{itemize}
\end{itemize}

\subsection{Language specific shortcuts:}
\label{sec:org8195ea0}
Replace \texttt{LANG} with \texttt{Julia}, \texttt{MatLab}, \texttt{Python}, or \texttt{R}, for the \href{https://julialang.org/}{Julia}, \href{https://www.mathworks.com/products/matlab.html}{MatLab}, \href{https://www.python.org/}{Python}, \href{https://www.r-project.org/}{R}
language, respectively, for the following commands.

\begin{itemize}
\item \texttt{\textbackslash{}runLANG[Arg1]\{Arg2\}\{Arg3\}[Arg4]} runs an external \texttt{LANG} code file.
\begin{itemize}
\item \texttt{Arg1} is optional and uses \href{https://pypi.org/project/talk2stat/}{talk2stat}'s \texttt{LANG} server by default.
\item \texttt{Arg2}, \texttt{Arg3}, and \texttt{Arg4} have the same effects as those of the basic command
\texttt{\textbackslash{}runExtCode}.
\end{itemize}
\item \texttt{\textbackslash{}runLANGIncOut[Arg1]\{Arg2\}[Arg3][Arg4][Arg5]} runs an external \texttt{LANG} code
file and embeds the output.
\begin{itemize}
\item \texttt{Arg1} is optional and uses \href{https://pypi.org/project/talk2stat/}{talk2stat}'s \texttt{LANG} server by default.
\item \texttt{Arg2}, \texttt{Arg3}, \texttt{Arg4}, and \texttt{Arg5} have the same effects as those of the command
\texttt{\textbackslash{}runCodeIncOut}.
\end{itemize}
\item \texttt{\textbackslash{}inlnLANG[Arg1]\{Arg2\}[Arg3][Arg4]} runs \texttt{LANG} source code (\texttt{Arg2}) and displays
the output in line.
\begin{itemize}
\item \texttt{Arg1} is optional and uses the \texttt{LANG} server by default.
\item \texttt{Arg2} is the \texttt{LANG} source code to run. If the \texttt{LANG} source code is wrapped
between "\texttt{```}" on both sides (as in the markdown grammar), then it will be
implemented directly; otherwise the code will be written to a file on the
disk and then be called.
\item \texttt{Arg3} and \texttt{Arg4} have the same effects as those of the basic command \texttt{\textbackslash{}inln}.
\end{itemize}

\item \texttt{\textbackslash{}runLANGChunk[Arg1]\{Arg2\}\{Arg3\}[Arg4][Arg5][Arg6]} runs a selected chunk from
an external \texttt{LANG} code file and embeds the output. 
\begin{itemize}
\item \texttt{Arg1} (optional ) uses \href{https://pypi.org/project/talk2stat/}{talk2stat}'s \texttt{LANG} server by default.
\item \texttt{Arg2} is the source file name.
\item \texttt{Arg3} is is the chunk identifier.
\item \texttt{Arg4} (optional) controls whether to run the code.
\item \texttt{Arg5} (optional) is the output file name.
\item \texttt{Arg6} (optional) controls the type of output with a default value \texttt{vbox}.
\end{itemize}
\end{itemize}

For example,
\begin{minted}[]{latex}
\runR{code/MontyHall_1.R}{montyhall-R1}
\runRIncOut{code/MontyHall_1.R}[][montyhall-R1]
\end{minted}

Specifically for \href{https://www.python.org/}{Python}, 

\begin{itemize}
\item \texttt{\textbackslash{}runPythonBatch[Arg1][Arg2]} runs an external \href{https://www.python.org/}{Python} code file in batch mode
(without a server running). Python (at least currently), unlike the other
languages we use, does not have an option to save and restore a session, which
means that once a Python session ends, the working environement (variable,
functions) is deleted. In order to allow a batch-mode in Python, we
implemented such capability. It requires the \href{https://pypi.org/project/dill/}{dill} module, which has to be
installed via \texttt{pip3 install dill}.
\begin{itemize}
\item \texttt{Arg1} is the \href{https://www.python.org/}{Python} source file name,
\item \texttt{Arg2} is the output file name.
\end{itemize}
\end{itemize}


\section{Revisions}
\label{sec:orge3b69ca}
\begin{itemize}
\item v2.4, January 17, 2025: 1. put file names within \$\$ for more robust
warnings; 2. use the python \href{https://pypi.org/project/advance-touch/}{advance-touch} package to create sub-folders for
chunk related commands.
\item v2.3, January 12, 2024: two bug fixes (1. removed an extra space after
\texttt{\textbackslash{}inlnX}; 2. when an underscore appeared inside R code, the tex file would
compile with an error).
\item v2.2, September 8, 2023: add \texttt{\textbackslash{}showChunk} basic command and \texttt{\textbackslash{}runLANGChunk}
commands for multiple languages.
\item v2.1, June 30, 2023: detokenize code which is passed to \texttt{\textbackslash{}inln}. This is
necessary when the code contains special latex characters like backslash.
These characters are escaped by latex when they are passed to macros as arguments,
which, in the case of \texttt{\textbackslash{}inln}, alters the correct syntax and causes execution errors.
Detokenizing ensures that the code is passed as-is to R/Julia/MatLab/Python.
We thank \href{https://github.com/kiryph}{kiryph} for reporting the issue.
\item v2.0, June 23, 2023: add \texttt{\textbackslash{}runCodeIncOut} command that runs an external code and
embeds the output and andd some language specific shortcuts; update
\texttt{\textbackslash{}runExtCode} so that if the output file does not exist, the \texttt{cache} option will
be overridden and the code will run; update \texttt{\textbackslash{}inln} so that its \texttt{Arg4} accept
\texttt{cache} or appending \texttt{.cache} to use the cached result. We thank \href{https://github.com/kiryph}{kiryph} for
suggesting these features.
\item v1.9, June 13, 2023: update \texttt{\textbackslash{}inln} command; the optional \texttt{Arg3} is the output
file name and the optional \texttt{Arg4} is the output type.
\item v1.8, January 18, 2023, add support to \href{https://ctan.org/pkg/listings?lang=en}{listings.}
\item v1.7, August 20, 2022: changed the tmp/ folder to generated/ in order to
conform with CTAN suggestions; renamed the troubleshooting file.
\item v1.6, August 10, 2022: stop only configured/running servers; a new
reducedspace option - some document classes put more space after the code box;
changed the default timeout of servers to 60 seconds; expanded the
troubleshooting document. New examples are now available on GitHub, including
how to collaborate with people who use Overleaf.
\item v1.5, July 23, 2022: Removed the utf8x option when loading inputenc due to a
conflict with hyperref.
\item v1.4, July 18, 2022: Fixed a bug in the cache mode.
\item v1.3, May 14, 2022: Removed the hard-coded minted options.
\item v1.2, May 3, 2022: Added python options (server and batch).
\item v1.1, April 17, 2021: Added a nohup option; improved error handling (missing
code files, zero bytes in output files.)
\end{itemize}

\section{Contributing}
\label{sec:orged22cfb}
We welcome your contributions to this package by opening issues on
GitHub and/or making a pull request. We also appreciate more example
documents written using \texttt{runcode}.


\textbf{Citing} \texttt{runcode}:
Haim Bar and HaiYing Wang (2021). \href{https://jds-online.org/journal/JDS/article/103/info}{Reproducible Science with \LaTeX{}},
\emph{Journal of Data Science} 2021; 19, no. 1, 111-125, DOI 10.6339/21-JDS998
\end{document}
